%%=============================================================================
%% Voorwoord
%%=============================================================================

\chapter*{\IfLanguageName{dutch}{Woord vooraf}{Preface}}%
\label{ch:voorwoord}

%% TODO:
%% Het voorwoord is het enige deel van de bachelorproef waar je vanuit je
%% eigen standpunt (``ik-vorm'') mag schrijven. Je kan hier bv. motiveren
%% waarom jij het onderwerp wil bespreken.
%% Vergeet ook niet te bedanken wie je geholpen/gesteund/... heeft

Als laatste loodjes van deze opleiding presenteer ik u graag mijn bachelorproef over hoe het KYC-proces bij de identiteitsverificatie van bedrijven op het Peppol-platform van Scrada geautomatiseerd kan worden. 
Deze werd opgesteld en geschreven tussen september 2025 en juni 2026.

Gedurende deze periode heb ik veel geleerd over het KYC-proces, de uitdagingen van identiteitsverificatie en de grenzen van automatisering. 
Als studente Full Stack Development, heb ik mijn technische vaardigheden kunnen uitbreiden en leren toepassen in een real-world context.
Na al de jaren kennis van de opleiding leek het me een interessante uitdaging om mijn huidige kennis over automatisering uit te breiden door onder andere gebruik te maken van AI-modellen en deze toe te passen.
In een tijd waar AI onvermijdelijk is leek het me vanzelfsprekend mezelf eens kans te geven dit zelf eens onder de loep te nemen. Hiermee ging nog steeds een software development aan vast,
dus ik heb zeker ook mijn software development skills hiermee kunnen aanscherpen. De keuze om dit onderwerp te behandelen is omdat ik al sinds mijn graduaatsopleiding in contact kwam met Scrada.
Ik had de kans gekregen om stage te doen bij hen, maar jammer genoeg lieten de omstandigheden dit niet toe. Uiteindelijk besloot ik om tijdens deze opleiding opnieuw contact op te nemen met Scrada voor mijn bachelorproef. 
Ik kreeg de kans om hun systeem te moderniseren door middel van automatisering, wat tegenwoordig van groot belang is. 
Deze onderzoeksvraag sloot meteen aan bij mijn persoonlijke motivatie binnen de IT-sector: het kunnen maken van een impact door iets dat zelf werd gemaakt of opgelost.

\medskip

Ik wil mijn bijzondere dank uiten aan mijn promotor Tom Antjon die altijd beschikbaar was wanneer nodig en voor het goede advies om me op weg te kunnen helpen.
Daarnaast wil ik natuurlijk ook mijn co-promotor Jeffrey Baeke van Scrada bedanken om deze kans te geven en me de juiste richtlijnen te geven om alles correct aan te pakken naar de wensen van Scrada.
Speciaal zou ik ook mijn dank willen uiten aan Ninho Decaesteker om me een duwtje in de rug te geven om het gebruik van de openbare KBO-databank kosteloos en correct te kunnen aanpakken.
Daarnaast zou ik ook graag mijn partner Yenthly Devolder en mijn broer Max Chandra willen bedanken, om me op weg te helpen met dingen die ik over het hoofd had gezien.
Tot slot wil ik natuurlijk ook mijn ouders bedanken die mij gedurende de hele opleiding me hebben gesteund en me tot dit punt hebben gebracht.

\medskip

Ik wens u veel leesplezier toe! \\Kira Hamers